@@M0@@ by (1.3.ii)).

Finally, observe that for @@M1@@, @@M2@@.

\sgaStableTarget{sga6.e05.thm.1-6}\textbf{Theorem 1.6.} \uline{The map from} @@M3@@ \uline{to} @@M4@@ \uline{which}, \uline{to a morphism of functors}, \uline{assigns its value on the element} @@M5@@ \uline{of} @@M6@@, \uline{is injective}; \uline{its image is the subgroup of} @@M7@@ \uline{formed by the elements} @@M8@@ \uline{such that} (\uline{with the notation of} 1.4. \uline{for the} @@M9@@):
@@M10@@

Let us first prove injectivity. By Lemma 1.2, it is enough to see that a morphism of functors @@M11@@ vanishing on @@M12@@ vanishes on @@M13@@. The commutative diagram:
\sgaStableTarget{sga6.e05.eq.1-6-1}@@M14@@
shows that for every @@M15@@, @@M16@@.

Now @@M17@@; hence @@M18@@, and consequently @@M19@@, for every @@M20@@; since @@M21@@ is separated and complete, @@M22@@.

Now let @@M23@@ be any element of @@M24@@, and put
